\documentclass[conference,a4paper]{IEEEtran}
\IEEEoverridecommandlockouts

% Packages retained from the manuscript; duplicate/conflicting imports removed.
\usepackage{cite}
\usepackage{float}
\usepackage{amsmath,amssymb,amsfonts}
\usepackage{mathtools}
\usepackage{algorithm}
\usepackage{algpseudocode}
\usepackage{graphicx}
\usepackage{textcomp}
\usepackage{xcolor}
\usepackage{booktabs}
\usepackage{tabularx}
\usepackage{array}
\usepackage{nicematrix}
\usepackage{balance}
\usepackage{enumitem}
\usepackage{tikz}
\usepackage{bm}
\usepackage{lipsum}
\usetikzlibrary{shapes.geometric}
\usepackage[hidelinks]{hyperref}

\algrenewcommand\algorithmicrequire{\textbf{Input:}}
\algrenewcommand\algorithmicensure{\textbf{Output:}}
\newtheorem{coro}{Corollary}
\newtheorem{theo}{Theorem}
\newtheorem{prop}{Proposition}
\newtheorem{remark}{Remark}
\newtheorem{example}{Example}
\newtheorem{lemma}{Lemma}
\newtheorem{defin}{Definition}
\hypersetup{breaklinks=true,citecolor=black,colorlinks=true,
            linkcolor=black,urlcolor=blue}
\setcounter{MaxMatrixCols}{30}

\title{TBD}
\author{\IEEEauthorblockN{TBD\IEEEauthorrefmark{1},
TBD\IEEEauthorrefmark{2}}
\IEEEauthorblockA{\IEEEauthorrefmark{1}XXX}
\IEEEauthorblockA{\IEEEauthorrefmark{2}Computer Communications Laboratory,
The University of Aizu, Aizuwakamatsu 965-8580, Japan}}

\begin{document}
\maketitle

\begin{abstract}
\lipsum[1]
\end{abstract}

\begin{IEEEkeywords}
Continuous-variable quantum key distribution (CV--QKD),
unmanned aerial vehicle (UAV), free-space optics (FSO),
high-altitude platform (HAP), probabilistic constellation shaping (PCS)
\end{IEEEkeywords}

\section{AUTOENCODER-BASED HAP--UAV CV-QKD SYSTEM}
\label{sec:autoencoder}

\subsection{Availability-Aware Joint UAV Position and Source Optimization}
\label{subsec:availability_aware_uav}

We consider a HAP--UAV MB-shaped DM-CV-QKD system under
moving-cloud conditions. The UAV flies at a fixed altitude
within a $1\,\mathrm{km}\times1\,\mathrm{km}$ area, discretized
into a $10\times10$ grid of $J=100$ candidate cell-center
positions $\mathcal Q=\{\mathbf q_1,\ldots,\mathbf q_J\}$.
Only one UAV is deployed; the grid points represent its
possible locations rather than separate ground stations.
The mission consists of $N_t$ time slots of duration
$\Delta t$, where $t_n=(n-1)\Delta t$.

To model moving clouds with low complexity, we consider a
single circular cloud at a fixed altitude $h_c$. Its
horizontal center evolves as
\begin{align}
\mathbf c_n
={}&\mathbf c_0+\mathbf v_c t_n,
\label{eq:availability_cloud_motion}
\end{align}
where $\mathbf c_0$ and $\mathbf v_c$ denote the initial center
and horizontal velocity. Let $\mathbf q_H$ be the HAP's
horizontal coordinates and let $h_H$ and $h_U$ be the HAP
and UAV altitudes, with $h_U<h_c<h_H$. The intersection
of the optical path from candidate position $j$ with the
cloud-altitude plane is
\begin{align}
\boldsymbol\ell_j
={}&\mathbf q_j+
\frac{h_c-h_U}{h_H-h_U}(\mathbf q_H-\mathbf q_j).
\label{eq:availability_cloud_intersection}
\end{align}
A cloud-blockage indicator and its additional power
transmittance are defined as
\begin{align}
b_{j,n}
={}&\mathbb I\!\left\{
\|\boldsymbol\ell_j-\mathbf c_n\|_2\leq R_c
\right\},
\label{eq:availability_blockage}
\\
T_{c,j,n}
={}&10^{-A_c b_{j,n}/10},
\label{eq:availability_cloud_transmittance}
\end{align}
where $R_c$ and $A_c$ are the cloud radius and additional
attenuation in dB, respectively. The original HAP--UAV
channel is augmented with $T_{c,j,n}$, without counting
cloud attenuation twice. The adopted QKD evaluator then
computes mutual information and Holevo information for
each updated channel state.

\noindent\textit{Remark 1.}
For candidate position $j$ and slot $n$, the asymptotic
secret key rate (SKR) is
\begin{align}
K_{j,n}(\nu_{\rm MB},V_A)
={}&\beta_{\rm rec}I_{AB,j,n}(\nu_{\rm MB},V_A)
\nonumber\\
&-\chi_{BE,j,n}(\nu_{\rm MB},V_A).
\label{eq:availability_skr}
\end{align}
Let $\mathcal C$ be the finite catalog of admissible,
exact-source parameter pairs $(\nu_{\rm MB},V_A)$, with
$\nu_{\rm MB}\geq0$ and $V_A>0$. At each position and
slot, the source parameters are chosen according to
\begin{align}
\mathcal P_1:\quad
\underset{(\nu_{\rm MB},V_A)\in\mathcal C}
{\operatorname{maximize}}\quad&
K_{j,n}(\nu_{\rm MB},V_A).
\label{eq:availability_source_optimization}
\end{align}
The resulting usable SKR is
\begin{align}
S_{j,n}
={}&\Big[\max_{(\nu_{\rm MB},V_A)\in\mathcal C}
K_{j,n}(\nu_{\rm MB},V_A)\Big]^+,
\label{eq:availability_usable_skr}
\end{align}
where $[x]^+=\max\{0,x\}$. The maximizing pair
$(\nu_{j,n}^{\star},V_{A,j,n}^{\star})$ is stored for each
state. This search does not require gradients of the SKR.

\subsubsection{QKD Availability and Stationary UAV Selection}
A QKD link is regarded as available when its optimized
SKR reaches a prescribed positive threshold $K_{\min}>0$.
Define
\begin{align}
a_{j,n}
={}&\mathbb I\!\left\{S_{j,n}\geq K_{\min}\right\},
\label{eq:availability_indicator}
\end{align}
so that $a_{j,n}\in\{0,1\}$ denotes QKD-link
availability at position $j$ in time slot $n$.
For a stationary UAV, its availability is
\begin{align}
A_j
={}&\frac{1}{N_t}\sum_{n=1}^{N_t}a_{j,n}.
\label{eq:stationary_availability}
\end{align}
The best fixed UAV location is determined by
\begin{align}
\mathcal P_2:\quad
\underset{j\in\{1,\ldots,J\}}{\operatorname{maximize}}
\quad&A_j.
\label{eq:availability_stationary_opt}
\end{align}
A prescribed reference position $j_{\rm ref}$ provides the
nonoptimized stationary baseline. The average SKR is also
reported to distinguish configurations with equal
availability but different key-generation rates.

\subsubsection{Mobility-Aware Availability Optimization}
For a moving UAV, let $j_n$ be its position index in slot
$n$ and define $d_{i,j}=\|\mathbf q_i-\mathbf q_j\|_2$.
The total flight distance and trajectory availability are
\begin{align}
D_{\rm fly}
={}&\sum_{n=2}^{N_t}d_{j_{n-1},j_n},
\label{eq:availability_flight_distance}
\\
A_{\rm UAV}
={}&\frac{1}{N_t}\sum_{n=1}^{N_t}a_{j_n,n}.
\label{eq:availability_trajectory}
\end{align}
To account for mobility expenditure without a detailed
propulsion model, we use the approximate incremental
movement-energy model $E_{\rm mov}=e_{\rm mov}D_{\rm fly}$,
where $e_{\rm mov}$ is a constant energy-per-distance
coefficient. With $E_{\rm ref}=e_{\rm mov}D_{\rm ref}$,
maximizing $A_{\rm UAV}-\lambda E_{\rm mov}/E_{\rm ref}$
is equivalent to solving
\begin{align}
\mathcal P_3:\quad
\underset{j_1,\ldots,j_{N_t}}{\operatorname{maximize}}
\quad&\frac{1}{N_t}\sum_{n=1}^{N_t}a_{j_n,n}
\nonumber\\
&-\frac{\lambda}{D_{\rm ref}}
\sum_{n=2}^{N_t}d_{j_{n-1},j_n}
\label{eq:availability_moving_opt}
\\
\text{subject to}\quad&j_n\in\{1,\ldots,J\},\ \forall n,
\label{eq:availability_grid_constraint}
\\
&d_{j_{n-1},j_n}\leq v_{\max}\Delta t,
\quad n\geq2,
\label{eq:availability_speed_constraint}
\end{align}
where $D_{\rm ref}>0$ is a reference flight distance,
$\lambda\geq0$ is a dimensionless tradeoff parameter, and
$v_{\max}$ is the maximum UAV speed. The initial position
is freely selectable before the mission, and deployment
costs are excluded from all baseline comparisons. Source
parameters may be updated once per time slot for all schemes.

\subsubsection{Low-Complexity Numerical Solution}
After solving $\mathcal P_1$, store $S_{j,n}$ and $a_{j,n}$
in $J\times N_t$ tables. Problem $\mathcal P_2$ is solved by
comparing $A_j$ for all $J$ positions. For $\mathcal P_3$,
let $F_n(j)$ be the maximum accumulated objective up to
slot $n$, conditioned on ending at location $j$. Set
$\mathcal N(j)=\{i:d_{i,j}\leq v_{\max}\Delta t\}$. Dynamic
programming yields
\begin{align}
F_1(j)
={}&\frac{a_{j,1}}{N_t},
\label{eq:availability_dp_init}
\\
F_n(j)
={}&\frac{a_{j,n}}{N_t}
+\max_{i\in\mathcal N(j)}\!\Bigg\{
F_{n-1}(i)
\nonumber\\
&\hspace{27mm}-\frac{\lambda}{D_{\rm ref}}d_{i,j}
\Bigg\},\quad n\geq2.
\label{eq:availability_dp_update}
\end{align}
The optimal final position is
$j_{N_t}^{\star}=\arg\max_jF_{N_t}(j)$. The optimal
trajectory on the candidate grid is recovered by
backtracking, as described in
Algorithm~\ref{alg:availability_joint_uav}.

\noindent\textit{Remark 2.}
The cloud evolution is assumed known for offline optimization.
The solution is optimal only over the considered grid and
source catalog. The linear movement-energy approximation
is not a calibrated UAV propulsion model. A strictly
positive $K_{\min}$ is required because setting $K_{\min}=0$
would classify all nonnegative $S_{j,n}$ values as available.
The adopted fading-security approximation must be validated
independently before interpreting numerical SKR results as
full fading-channel security guarantees.

\begin{algorithm}[t]
\caption{Availability-Aware Joint Source and UAV Position Optimization}
\label{alg:availability_joint_uav}
\begin{algorithmic}[1]
\Require $\mathcal Q$, $\mathcal C$, cloud parameters,
$N_t$, $\Delta t$, $v_{\max}$, $K_{\min}$,
$\lambda$, $D_{\rm ref}$.
\For{$n=1,\ldots,N_t$}
  \For{$j=1,\ldots,J$}
    \State Update the cloud-dependent channel.
    \State Search $\mathcal C$ and store $S_{j,n}$,
    $(\nu_{j,n}^{\star},V_{A,j,n}^{\star})$.
    \State Set $a_{j,n}\gets\mathbb I\{S_{j,n}\geq K_{\min}\}$.
  \EndFor
\EndFor
\State Solve $\mathcal P_2$ by comparing all $A_j$.
\State Set $F_1(j)\gets a_{j,1}/N_t$ for all $j$.
\For{$n=2,\ldots,N_t$}
  \For{$j=1,\ldots,J$}
    \State Update $F_n(j)$ using
    \eqref{eq:availability_dp_update}.
    \State Save the optimal predecessor of $j$.
  \EndFor
\EndFor
\State Select $\arg\max_jF_{N_t}(j)$ and backtrack.
\Ensure Stationary location, UAV trajectory,
source parameters, QKD availability, average SKR,
and flight distance.
\end{algorithmic}
\end{algorithm}

\end{document}
